% ---------------------------------------------------------------------------
% Revised abstract. Drop-in replacement for the \begin{abstract}...\end{abstract}
% block in shift-surv-paper-PATCHED.tex.
%
% Incorporates the review comments, plus one correction that came out of
% checking the code:
%
%   "best adaptation method" is chosen PER TRANSFER as the one with the largest
%   observed gap on the target (meta_analysis.py:164, pool_best_per_transfer).
%   That selection sees target labels, so it is a winner's-curse upper bound --
%   deliberately optimistic. Stating it makes the negative result stronger, not
%   weaker: not even a retrospectively chosen best method beats ERM. Left
%   unstated, a reviewer who finds it reads it as a concealed optimism.
%
% Terminology: "label-prior shift" is kept as the scenario name and tied to the
% label-shift literature at first mention, because "Label-prior" is baked into
% seven figure images. Switching to "label shift" everywhere means regenerating
% all of them.
%
% Length: ~270 words, down from ~400. IEEE conference abstracts run 150-250.
% ---------------------------------------------------------------------------

\begin{abstract}
Clinical risk models often lose performance when applied at centers other than
the one where they were developed. Unsupervised domain adaptation (DA) and
domain generalization (DG) are widely proposed remedies, but real-world data
rarely identify which shift mechanism caused the degradation, so method failure
cannot be separated from an intrinsically unsolvable transfer. We introduce
SHIFT-SURV, a controlled benchmark for fixed-horizon survival-risk prediction
that independently varies covariate, measurement, spurious-association,
baseline-hazard, label-prior, and concept shift against a known invariant
oracle. Continuous event times are generated from a latent-factor Weibull
proportional-hazards model, with invariant and spurious factors mapped to
tabular, biomarker-like covariates and dichotomized at a six-month horizon.

Across per-center sample sizes of 250--2{,}000, CORAL, DANN, CDAN, IRM, and
GroupDRO do not outperform empirical risk minimization (ERM) on target AUROC,
including where the oracle shows substantial recoverable signal: under
measurement shift the ERM-to-oracle AUROC gap reaches $0.21$ and remains
essentially unrecovered at the largest sample size. Mechanistic analysis shows
why. CORAL halves the source--target covariance discrepancy while leaving the
angle between the learned predictive direction and the target-optimal direction
essentially unchanged.

Discrimination and calibration deteriorate independently. Unsupervised
recalibration improves calibration under pure label-prior shift, increasingly so
with the amount of unlabeled target data, and worsens it when its assumptions
are violated. Across center transfers in I3LUNG and a public breast-cancer
registry, a random-effects meta-analysis finds no clinically meaningful gain
over ERM for the best adaptation method, selected per transfer by the largest
observed target gap and therefore an optimistic upper bound: its 95\% confidence
interval excludes an absolute AUROC improvement of $0.02$. SHIFT-SURV separates
irreducible transport difficulty from practical method failure and identifies
when unsupervised transport methods fail in survival-risk prediction.
\end{abstract}
